\chapter{Introduction}

This survey is based on a course presented by the first author at the Research School in Marseille, March 6-10, 2017. 
The school was organized by the chair holders, Kostya Khanin and Senya Shlosman, of the Jean-Morlet chair 2017 of CIRM,
\begin{quotation}
Random Structures in Statistical Mechanics and Mathematical Physics.
\end{quotation}
The model is a space-time continuous directed polymer in random environment. 
In this regard, it is one of the most basic such model and it plays a fundamental
role. % by its simplicity. 
Directed polymers are described by random paths, which are influenced by randomly located impurities which may be attractive or repellent.
Such models have been widely considered in statistical physics, disordered systems and stochastic processes. 
%; see   \cite{CStFlour} for a recent review in the discrete setup.

\medskip


As an informal definition we model the polymer by a random path $\bx=(\bx(t); t \geq 0)$ taking values in $\rd$ and interacting with time-space Poisson points $(t_i, x_i)$ called environment. The path sees such a point if at time $t_i$ it is located within a fixed distance $r$ from $x_i$. Denoting by $\#_t(\bx)= \sum_{i: t_i\leq t}  {\bf 1}_{|\bx(t_i)-x_i|\leq r}$ the number of  Poisson points seen by the path $\bx$ up to time $t$,
the model with time horizon $t$ at inverse-temperature parameter $\b$ is associated to the Hamiltonian 
$$
-\frac 12 \int_0^t |\dot \bx (s)|^2 ds + \b \#_t(\bx)\;.
$$


In this model where the path is Brownian and the medium is Poissonian,  we benefit from nice formulas and strong tools from stochastic calculus for
Gaussian or Poisson measure and martingale techniques. 
\medskip


The notes are essentially based on references \cite{CY05, CYkokyuroku, CYBMPO2,CoscoIntermediate}, gathering and unifying the matter scattered in these references, and containing novel contributions and perspectives as emphasized below. It also parallels the book \cite{CStFlour} which deals similar models in the  discrete framework, and we warn the reader of the existence of many results available for one 
%of these particular models . 
particular model but  not for the others. We do not reproduce all details or computations, but we rather try to give the general picture and the essential arguments. 
\medskip

Let us mention the main highlights in this survey and also the new results:
 \begin{enumerate}
 \item We establish in section \ref{ch:3} a fine continuity estimate under spatial shifts for the limit of the martingale. This is achieved by a smart use of mirror coupling. 
\item Section \ref{ch:dir} contains a nice original account on directional free energy. We develop a full approach of
disorder strength based on directional free energy.
\item In section \ref{ch:CM} we develop an original approach to diffusivity at weak disorder, based on Camer\'on-Martin transformation (see theorem \ref{prop:diff}).
\item Section \ref{sec:IR} is dedicated to the intermediate disorder regime and KPZ equation. We give a synthetic account with all the central ideas. 
\end{enumerate}

The detailed matter and the organization appear most clearly in the table of contents,
which is a useful source to follow the line all through the notes. 

